% Source: https://artofproblemsolving.com/community/c3h1370809_handbook_problems
% Stable identity: aops-topic-1370809-q02
If the radius of circle $A$ has grown by $t$, the radii are $4+t$ and $6+2t$. The circles are externally tangent when the sum of the radii equals the distance between the centres, $40\sqrt2$:
\[
(4+t)+(6+2t)=40\sqrt2\iff t=\dfrac{40\sqrt2-10}{3},
\]
so the radius of circle $A$ is then $4+t=\dfrac{40\sqrt2+2}{3}$.

\begin{asy}
usepackage("amsmath");
size(5.5cm);
defaultpen(fontsize(10pt));
real s2=sqrt(2), t=(40*s2-10)/3, ra=4+t, rb=6+2t;
pair A=(0,0), B=(40,40), P=ra*dir(45);
draw(circle(A,4),gray+dashed); draw(circle(B,6),gray+dashed);
draw(circle(A,ra),blue); draw(circle(B,rb),red);
draw(A--B,dashed);
dot(A);dot(B);dot(P);
draw(P--(P.x,0),gray+dotted); draw((-ra-2,0)--(B.x+rb+2,0),gray);
label("$A$",A+5.2*dir(225),SW); label("$B$",B+(6.2,0),E); label("$P$",P,1.4*W);
label("$x$",(P.x,0),S);
label("\mbox{}",(min(currentpicture,user=true)+max(currentpicture,user=true))/2,black); // keeps the text after the figure black when LaTeX typesets the labels inline
\end{asy}

The point of tangency $P$ lies on segment $AB$, which lies along the line $y=x$, at distance $\dfrac{40\sqrt2+2}{3}$ from the origin. So
\[
x=\dfrac{40\sqrt2+2}{3}\cdot\dfrac{1}{\sqrt2}=\boxed{\dfrac{40+\sqrt2}{3}}\approx13.81 .
\]
